\documentclass[11pt]{amsart}
\usepackage{geometry}                % See geometry.pdf to learn the layout options. There are lots.
\geometry{letterpaper}                   % ... or a4paper or a5paper or ... 
%\geometry{landscape}                % Activate for for rotated page geometry
%\usepackage[parfill]{parskip}    % Activate to begin paragraphs with an empty line rather than an indent
\usepackage{graphicx}
\usepackage{amssymb}
\usepackage{epstopdf}
\usepackage{multirow}
\DeclareGraphicsRule{.tif}{png}{.png}{`convert #1 `dirname #1`/`basename #1 .tif`.png}

\title{Problema 5.3}
\author{}
%\date{}                                           % Activate to display a given date or no date

\begin{document}
\maketitle
%\section{}
%\subsection{}
En una universidad almacenan información relativa a sus profesores en un arreglo unidimensional. Cada elemento del arreglo es un registro que contiene la siguiente información respecto a cada profesor:

\begin{itemize}
  \item Número de empleado.
  \item Nombre del profesor.
  \item Departamento al que pertenece.
  \item Grado académico.
  \item Nacionalidad.
  \item Salario (Se almacena en forma mensual lo que cobra el profesor en un arreglo unidimensional).
\end{itemize}

Construya un diagrama de flujo que pueda proporcionar la siguiente información:

\begin{enumerate}
  \item El número, nombre y nacionalidad del profesor que más ganó el año anterior.
  \item El monto total pagado en el año a los profesores del departamento X.
  \item El monto total pagado a los profesores extranjeros (nacionalidad distinta a México).
  \item El nombre del profesor del departamento Y que más ganó en el año anterior.
\end{enumerate}

Dato: PROFESORES [1..N] 1 < N < 200

Donde:

PROFESORES es un arreglo unidimensional de registros.


\begin{itemize}
  \item NUM: Variable de tipo entero que representa el número del empleado.
  \item NOM: Variable de tipo cadena de caracteres que indica el nombre del profesor.
  \item DEP: Variable de tipo cadena de caracteres que representa el departamento en que está adscrito el profesor.
  \item GRADO: Variable de tipo cadena de caracteres que expresa el máximo grado académico del profesor.
  \item NAC: Variable de tipo cadena de caracteres que representa la nacionalidad del profesor.
  \item SALARIO: Arreglo unidimensional de tipo real que almacena en forma mensual lo que cobró el profesor.
\end{itemize}

\end{document}
